\begin{table}[htbp]
\centering
\caption{PHFSI Weighting Scheme Comparison: Equal vs PCA}
\label{tab:pca_comparison}
\footnotesize
\begin{threeparttable}
\begin{tabular}{lccccc}
\toprule
Weighting Scheme & Mean & SD & Min & Max & Correlation \\
\midrule
Equal weights & 0.465 & 0.196 & 0.137 & 1.000 & — \\
PCA weights & 0.659 & 0.100 & 0.055 & 0.947 & 0.683 \\
\midrule
\multicolumn{6}{l}{\textit{Spearman rank correlation: } ρ = 0.790} \\
\multicolumn{6}{l}{\textit{Mean absolute difference: } 0.224} \\
\bottomrule
\end{tabular}
\begin{tablenotes}
\small
\item \textit{Notes:} Comparison of PHFSI calculated using equal weights (0.20 each for 5 components) vs PCA-derived weights from first principal component (variance explained: 74.7\%). PCA weights: OSSR=0.258, LRR=0.260, TLR=0.235, SPI=0.246. High correlation (r=0.683) indicates weighting scheme does not materially affect index rankings. N=435 hospital-year observations with complete component data.
\end{tablenotes}
\end{threeparttable}
\end{table}